%!TEX root=../main.tex
\section{Discussion and Conclusion}

  We have observed cases in which calibration and relaxed Equalized Odds
  are compatible and cases where they
  are not. When it is feasible, the penalty of equalizing cost is amplified if
  the base rates between groups differ significantly. This is expected,
  as base rate differences are what give rise to cost-disparity in the calibrated setting.
  Seeking equality with respect to a single error rate (e.g. false-negatives, as in the income
  prediction experiment) will necessarily increase disparity with respect to the
  other error. This may be tolerable (in the income prediction case, some
  employees will end up over-paid) but could also be highly problematic (e.g.
  in criminal justice settings).
%
  % metric, while simultaneously increasing error.
  % If we choose to equalize one of the error rates while maintaining calibration,
  % then we will necessarily increase disparity with respect to the other error
  % metric, while simultaneously increasing error.
%
  Finally, we have observed that the calibrated relaxation is infeasible when the best (discriminatory)
  classifiers are not far from the trivial classifiers (leaving little room for interpolation). In such settings,
  we see that calibration is completely incompatible with an equalized error
  constraint.

In summary, we conclude that maintaining cost parity \emph{and} calibration is
desirable yet often difficult in practice. Although we provide an algorithm to
effectively find the unique feasible solution to both constraints, it is
inherently based on randomly exchanging the predictions of the better
classifier with the trivial base rate. Even if fairness is reached in
expectation, for an individual case, it may be hard  to accept that
occasionally consequential decisions are made by randomly withholding
predictive information, irrespective of a particular person's feature representation. In this
paper we argue that, as long as calibration is required, no lower-error solution can be achieved.


  % The only setting in which the calibrated framework yields reasonable results is
  % when we are concerned with matching a weighed combination of error rates.
  % In this setting, we are able to maintain calibration, while also introducing
  % less error than if we explicitly matched both error rates. However, this relaxation
  % only makes sense in certain real-world scenarios. In the criminal recidivism
  % example, it would be possible to achieve this notion of non-discrimination by
  % disproportionately assigning false-positive errors to one group while assigning
  % false-negative errors to the other (as is the case now). While both errors
  % are unfavorable to society at large, it is fair to argue that false-negative errors
  % are favorable to the defendants while false-positive errors are especially
  % detrimental. Therefore, the weighted cost notion of non-discrimination is
  % not applicable to all settings, and in some cases could incur new types of
  % discrimination.


% An increasing amount of attention is being devoted to the issue of fairness in machine learning.  In this paper we specifically investigate settings in which calibration and
% disparity of errors are important.
% We introduce \methodname{}, a relaxation of Equalized Odds that maintains
% classifier calibration. To achieve this notion of non-discrimination, we
% introduce a post-processing method which is not only
% feasible but optimal, and we also show fundamental limitations on how many constraints
% calibrated classifiers can satisfy. Finally, we empirically evaluate the trade-offs between
% this calibrated notion of non-discrimination and the original Equalized Odds framework.
% It is important to note that the intent of our framework is not to resolve the
% tradeoffs inherent to fair classification; we instead seek to make explicit a
% space of possible solutions. This gives a practitioner the opportunity to
% choose which framework is most applicable to the scenario he or she faces.

% While our method -- and indeed any method that seeks to achieve similar
% constraints -- necessarily trades off performance to reduce disparity, we note
% that this tradeoff will incentivize the designer of a classifier to
% reduce errors for the group that suffers higher cost, as the overall
% performance of the system is limited by the performance on the worse-off group.
% This might motivate the collection of richer features or more
% data for the group with higher cost, leading to lower cost classifiers.
%
% Our work leads to several open questions related to fairness in
% classification. While our results hold for a broad class of disparity
% constraints, it would be interesting to further characterize
% constrained calibrated classification.
% Furthermore, we propose the more general problem of investigating tradeoffs
% between calibration and other constraints when they are not simultaneously
% satisfiable.
